\cotefile{156-8}
\scanpage{18}
\mspage{9}
% archive sheet number: 18
% lines 1--2
\[
\begin{aligned}
\text{b/}\quad &\unsure{\operatorname{envelopp}|\operatorname{envelop}}(F)\unsure{\ll|\subset\subset}\operatorname{envelopp}G\Longrightarrow F\ll G\\
\text{c/}\quad &\operatorname{envelopp}(F)\unsure{\gtrless|\bowtie}\operatorname{envelopp}(G)\Longrightarrow F\gtrless G.
\end{aligned}
\]
% lines 3--7
C'est quand ces \deleted{\unsure{conditions}} implications sont satisfaites \inserted{(en plus de la condition a) ci-dessus)} (et il suffit de le poser pour $F=X$, $\deleted{\unsure{F}}G=Y$, $X,Y\unsure{\in M}$ \inserted{(spécimens)}) qu'on dit que le \unsure{magasin} est ``ensembliste''.

\deleted{\unsure{Il est donc donné par un ensemble de}}

\deleted{Une telle donnée \unsure{équivaut|équival}}
% lines 8--14
\uline{Proposition}\inserted{-- \uline{Scholie}}. La donnée d'un \unsure{magasin} ensembliste, \unsure{à isom. près} revient à la donnée d'un ens. $L$ (l'ens des ``lieux'' ou ``points''), et d'une partie
\[
M\subset\operatorname{Figel}(L),
\]
satisfaisant les conditions
% lines 15--20
\uline{\unsure{Mag ens|Magas}} 1\quad $\forall\deleted{\unsure{\mathscr{F}}}F\in M$, et tout $\deleted{\unsure{x}}X\in\deleted{\unsure{\mathscr{F}}}F$, l'ens.
\[
\deleted{\unsure{\mathscr{F}}}F_{\unsure{\leq X|\subset X|X}}=(Y\in F\mid Y\subset X)\subset F
\]
est aussi dans $M$ ($M$ stable par \unsure{passage aux sous-multistrates}).

\uline{\unsure{Mag ens|Magas}} 2\quad $\forall x\in L$, $\{\{x\}\}\in M$.
% slanted
\margin{!!! Donc on a une \unsure{espèce de structure} $M\in\unsure{\mathscr{P}|\mathcal{P}}\mathscr{P}(L)$, avec \unsure{axiomes} \uline{\unsure{Mag ens|Magas}} 0 : $\forall X\in M$ est une figure ensembliste élémentaire, -- plus les deux \unsure{autres}, à peine plus \unsure{compliqué} que la définition d'un \unsure{espace} topologique !}
% lines 21--28
\unsure{Ça} c'est donc pas mal dégonflé !!!

En fait, à ce stade de l'axiomatique \unsure{``naïviste : là''|``naïviste'' là}, le point de vue ``atlas'' des figures nous avait plutôt compliqué la vie. Ici il vaut décidément mieux partir du \unsure{magasin}, et y voir aussi clair que possible -- quitte à faire ensuite une \unsure{place} \unsure{convenable|amplement} \unsure{aux|sur les} atlas plus
